\section{Introduction}
\label{sec:intro}

The Constraint Satisfaction Problem (CSP) framework is foundational to modeling complex combinatorial search and placement problems in computer science and artificial intelligence \citep{russell2010artificial}. Within this domain, the N-Queens problem---the challenge of placing $N$ non-attacking queens on an $N \times N$ chessboard---serves as a canonical benchmark for evaluating the efficiency and scalability of exact search algorithms \citep{bell2009solving}. As the problem size $N$ increases, the state space expands exponentially, rendering naive enumeration computationally intractable and necessitating sophisticated algorithmic paradigms.

Currently, exact solvers rely on three primary modeling paradigms: Boolean Satisfiability (SAT), Constraint Programming (CP), and Mixed Integer Programming (MIP) \citep{gomes2008satisfiability, rossi2006handbook, wolsey2020integer}. In the Boolean SAT paradigm, high-level structural constraints must be transformed into Conjunctive Normal Form (CNF) via variable-heavy At-Most-One (AMO) encodings. Different encoding strategies---such as Pairwise, Binary, Sequential Counter, Commander, and Product---drastically alter the size and structure of the resulting CNF formula \citep{sinz2005towards, frisch2005sat}. Conversely, native Constraint Programming implementations maintain the problem's global structural semantics by directly applying constraints over integer variables \citep{regin1994filtering}, while Mixed Integer Programming leverages linear inequality relaxations over binary decision bounds.

A recurring assumption in SAT modeling is that structural complexity---specifically the total number of generated clauses or auxiliary variables---correlates directly with the time required by Conflict-Driven Clause Learning (CDCL) solvers to find a solution. However, real-world benchmarking often reveals non-linear scaling behaviors heavily influenced by the solver's internal heuristics \citep{marques2009conflict}. Furthermore, while algorithmic complexity theoretically bounds performance, the practical runtime characteristics across the SAT, CP, and MIP paradigms on the exact same problem formulations remain sparsely compared in unified experimental setups.

The objective of this research is to execute a strictly controlled empirical comparison of these exact solving techniques on the N-Queens problem up to $N=100$. Specifically, we address the following research questions:
1. How do different AMO encodings affect the structural complexity of SAT formulations?
2. How does the choice of SAT encoding impact solving time and completion reliability under a uniform CDCL configuration?
3. How do general-purpose Boolean reasoning frameworks compare against specialized native CP and MIP solvers?
4. What are the practical scalability limitations of these paradigms?

To answer these questions, we constructed an automated, subprocess-isolated benchmarking framework. We evaluate five SAT encodings solved via PySAT Glucose3, against OR-Tools CP-SAT, IBM CP Optimizer, Gurobi MIP, and IBM CPLEX MIP. 

The primary contributions of this paper are:
\begin{itemize}
\item A comprehensive, reproducible empirical dataset reporting 540 scheduled experimental records (495 at $N \le 100$ and 45 at $N=200$) across 9 distinct exact solving methodologies, including 450 successfully validated solutions and explicitly documented timeout, licensing, and resource-policy outcomes.
\item Evidence that structural compactness (clause count minimization) is a poor predictor of CDCL resolution time on placement constraints.
\item The identification of severe non-monotonic search pathologies in specific structural encodings (e.g., the SAT-Product encoding) under standard phase-selection heuristics.
\item A direct comparative baseline demonstrating the scalability advantages of native Constraint Programming ($O(N)$ formulation) while highlighting specific SAT encodings that remain highly competitive at scale.
\end{itemize}
